% paper/sections/11d_coupling.tex
%
% §11.4  二相流時間積分：CLS↔PPE 結合検証
% ═══════════════════════════════════════════════════════════════════════════════
\subsection{二相流時間積分：CLS↔PPE 結合検証}
\label{sec:advection_pressure_coupling}
% ═══════════════════════════════════════════════════════════════════════════════

§\ref{sec:force_balance}--§\ref{sec:ns_time_accuracy} では単相 NS ソルバの整合性を確認した．
§~\ref{sec:bf_static_droplet} の Laplace 圧テストは変密度 PPE が\textbf{1ステップ}で正しく機能することを
示したが，圧力場はほぼゼロから始まる初期条件であり，Hermite 延長は実質的に no-op であった．
本節では，\textbf{多ステップ時間積分}で非零圧力場における Hermite 延長・CLS 移流・変密度 PPE の
結合が正しく機能するかを検証する．

\noindent\textbf{目的：}
完全な二相流タイムステップループ
（CLS 界面移流 → 物性値更新 → 変密度 PPE → Hermite 延長 → Corrector）が
物理的整合性を維持することの確認．特に以下を検証する：
\begin{enumerate}[noitemsep]
  \item 移動界面の存在下で $\bnabla\cdot\bu = 0$ が維持されるか（CLS↔PPE 結合）
  \item 非零圧力場での Hermite 延長が動的問題で機能するか（§~\ref{sec:verify_field_extension} 静的テストの拡張）
  \item 低密度比での CSF 体積力と変密度 PPE の整合性
\end{enumerate}

\noindent\textbf{前提となる検証済み結果（§~\ref{sec:numerical_integrity}章）：}
\begin{itemize}[noitemsep]
  \item CLS 質量保存誤差：§~\ref{sec:cls_conservation_bench}（$N{=}256$ で $\Ord{10^{-5}}$）
  \item 界面越し Hermite 延長精度：§~\ref{sec:verify_field_extension}（$N{=}128$ で相対誤差 $2.2 \times 10^{-3}$）
  \item 変密度 PPE（滑らか密度場）：§~\ref{sec:verify_ppe_varrho}（$\rho_{\max}/\rho_{\min}{\approx}999$ で $\Ord{h^{5.8}}$）
  \item AB2 時間精度：§~\ref{sec:time_accuracy_table}（$\Ord{\Delta t^2}$）
\end{itemize}

\subsubsection{セットアップ}

\begin{itemize}[noitemsep]
  \item \textbf{計算領域}：$[0,1]^2$，壁面境界条件（$\bu = \bm{0}$）
  \item \textbf{初期条件}：円形液滴（$R = 0.25$，中心 $(0.5, 0.5)$），初速 $\bu_0 = (1, 0)^T$（一様移流）
  \item \textbf{密度比}：$\rho_l/\rho_g = 2$，$5$（合格期待），$10$（警告期待）
  \item \textbf{パラメータ}：$We = 10$，$Re = 100$，$Fr^{-1} = 0$（重力なし）
  \item \textbf{格子}：$N = 32, 64, 128$
  \item \textbf{積分時間}：$T = 0.5$（CFL $= 0.5$ 基準の $\Delta t$）
  \item \textbf{評価指標}：$\|\bnabla\cdot\bu\|_\infty$ の時系列，質量誤差の時系列，$\Delta p$ 相対誤差
\end{itemize}

\subsubsection{合否基準}

\begin{table}[htbp]
\centering
\caption{二相流時間積分テストの合否基準}
\label{tab:twophase_ti_criteria}
\renewcommand{\arraystretch}{1.3}
\begin{tabular}{clll}
\toprule
$\rho_l/\rho_g$ & 判定 & $\|\bnabla\cdot\bu\|_\infty$ & 質量誤差 \\
\midrule
$2$ & PASS 期待 & $< 10^{-8}$ & $< 10^{-4}$ \\
$5$ & PASS 期待 & $< 10^{-8}$ & $< 10^{-4}$ \\
$10$ & WARN（精度低下許容） & $< 10^{-4}$ & $< 10^{-3}$ \\
\bottomrule
\end{tabular}
\end{table}

\noindent\textbf{§~\ref{sec:bf_static_droplet} との違い：}
§~\ref{sec:bf_static_droplet}（Laplace 圧）は単ステップ投影で圧力初期値ゼロから始めるため
Hermite 延長が no-op となる特殊条件であった．
本テストは多ステップ積分で前ステップの非零圧力場を Hermite 延長した後に
$-\bnabla p^n$ を Predictor に渡す通常運転条件の初の検証である．
$\rho_l/\rho_g = 10$ での精度劣化の兆候は §~\ref{sec:interface_crossing} の
破綻限界定量化テストへの橋渡しとなる．

\subsubsection{予測結果と技術的根拠}

§~\ref{sec:verify_ppe_varrho} で確認した通り，
変密度 PPE 欠陥補正は\textbf{滑らかな密度場}（$C^\infty$）では $\rho_{\max}/\rho_{\min} \approx 999$
においても $\Ord{h^{5.8}}$ を維持する（§~\ref{sec:varrho_ppe_limitation}）．
一方，界面型（smoothed Heaviside）密度場では $1/\rho$ のジャンプにより
$\rho_l/\rho_g \geq 10$ で発散することが予測される．
$\rho_l/\rho_g \leq 5$ は遷移層幅 $\varepsilon = 1.5h$ 内のジャンプが
$5/1 = 5$ 倍程度にとどまり，欠陥補正の収束が維持される可能性が高い．

\begin{tcolorbox}[mybox, title={実験状況：計画}]
本テストは実験スクリプト（\texttt{experiment/ch11/exp11\_4\_twophase\_ti.py}）の実装後に実施予定である．
上記の予測結果は §~\ref{sec:verify_ppe_varrho} および §~\ref{sec:varrho_ppe_limitation} の
コンポーネント検証結果からの外挿に基づく．
\end{tcolorbox}
